% Part: first-order-logic
% Chapter: beyond
% Section: other-logics

\documentclass[../../../include/open-logic-section]{subfiles}
% 这是一本关于数理逻辑的书，由 Open Logic Project 编写，主要面向哲学专业的学生，同时也兼顾数学与计算机科学专业的学生。

\begin{document}

\olfileid{fol}{byd}{oth}

\olsection{其他逻辑系统}

到现在你可能已经发现，设计一种新逻辑并非难事。
你也可以创造自己的语法，构造一个演绎系统，并为它配上相应的语义。
如果你希望推导系统关于语义是完备的，那或许需要一点巧思；
并且你可能得花些力气，才能让世人相信你的逻辑确实有趣。
但作为回报，你可以享受数小时纯粹而美好的乐趣，去探索你的逻辑所具有的数学与计算特性。

近几十年来，形式逻辑迎来了名副其实的爆发式发展。
模糊逻辑旨在为关于模糊性质的推理建模；
概率逻辑旨在为关于不确定性的推理建模；
默认逻辑和非单调逻辑旨在为可废止的推理形式建模，
也就是说，这些“合理的”推断在面临新信息时可能会被推翻。
还有认知逻辑，旨在为关于知识的推理建模；
因果逻辑，旨在为关于因果关系的推理建模；
甚至还有“道义”逻辑，旨在为关于道德与伦理义务的推理建模。
引入这些系统的主要动机可能是哲学、数学或计算方面的；
因此，你可能会在数理逻辑、哲学逻辑、人工智能、认知科学或其他领域的范畴下接触到这些逻辑系统。

这样的例子不胜枚举，可能性似乎无穷无尽。
我们或许永远无法实现 Leibniz将人类理性全部化归为计算的梦想——但这无法阻止我们去尝试。

\end{document}